%
% 22-partialbruchzerlegung.tex -- Partialbruchzerlegung
%
% (c) 2026 Prof Dr Andreas Müller
%
\subsection{Partialbruchzerlegung
\label{buch:analytisch:rational:subsection:partialbruch}}
Der Definitionsbereich einer rationalen Funktion $f(z)=p(z)/q(z)$
wird durch die Nullstellen des Nenners bestimmt.
Ein Integral entlang eines Weges, der eine Nullstelle des Nenners
umläuft, dann ist das Integral durch Funktionswerte und 
Ableitung in dieser Nullstelle bestimmt.
Es ist daher naheliegend, die rationale Funktion $f(z)$ in
einzelne Summanden zu zerlegen, welche jeweils nur eine einzige
Nullstelle im Nenner haben.

%
% Divisionsrest abspalten
%
\subsubsection{Divisionsrest abspalten}
In einem ersten Schritt kann man den Polynomdivisionsalgorithmus
durchführen und den Zähler als $p(z) = a(z)q(z) + r(z)$ mit
Polynomen $a(z)$ und $r(z)$ schreiben.
\index{Polynomdivision}%
Der Grad $\deg r(z)$ des Restes $r(z)$ muss kleiner sein als
der Grad von $q(z)$: $\deg r(z) <\deg q(z)$.
Damit kann die rationale Funktion in zwei Terme
\[
f(z)
=
\frac{a(z)q(z)+r(z)}{q(z)}
=
a(z)
+
\frac{r(z)}{q(z)}
\]
aufgeteilt werden.
Das Polynom $a(z)$ sind holomorph in ganz $\mathbb{C}$ und trägt
daher nicht zu Wegintegralen über geschlossene Wege bei.
Wir wollen daher im Folgenden annehmen, dass $\deg p(z)<\deg q(z)$.

%
% Nur eine Nullstelle des Nenners
%
\subsubsection{Nur eine Nullstelle des Nenners}
Wir betrachten zunächst den Spezialfall einer einzigen Nullstelle
$ a\in\mathbb{C}$ des Nenners.
\index{Nullstelle}%
Der Nenner muss somit die Form $q(z) = (z- a)^n$ haben.
Da der Bruch $p(z)/q(z)$ bereits gekürzt ist, muss $p( a)\ne 0$
sein.
Da $p(z)-p( a)=0$ ist, muss $p(z)-p( a)$ durch $z- a$
teilbar sein, es gibt also ein Polynom $p^1(z)$ mit der Eigenschaft
\[
p^1(z)(z- a) + p( a) = p(z).
\]
Der Grad von $p^(z)$ ist $\deg p^1(z) = \deg p(z)-1 < \deg p(z)$.
Durch wiederholte Anwendung dieses Prozesses entstehen Polynome
$p^2(z),\dots,p^m(z)$, mit 
\begin{equation}
\begin{aligned}
p^2(z)(z- a) + p^1( a)     &= p^1(z)           \\
p^3(z)(z- a) + p^2( a)     &= p^2(z)           \\[-2pt]
                                   &\phantom{a} \vdots \\
p^m(z)(z- a) + p^{m-1}( a) &= p^{m-1}(z).
\end{aligned}
\label{buch:analytisch:rational:partial:eqn:pk}
\end{equation}
Zur Vereinheitlichung der Notation schreiben wir zudem $p^0(z)=p(z)$.
Kürzen wir $A_1=p( a)$, $A_2=p^1( a)$ oder allgemein
$A_k = p^{k-1}(z)$ ab, dann werden die
Gleichungen~\eqref{buch:analytisch:rational:partial:eqn:pk}
zu
\[
\begin{aligned}
p^1(z)(z- a) + A_1 &= p^0(z)           \\
p^2(z)(z- a) + A_2 &= p^1(z)           \\
p^3(z)(z- a) + A_3 &= p^2(z)           \\[-2pt]
                       &\phantom{a} \vdots \\
p^m(z)(z- a) + A_m &= p^{m-1}(z).
\end{aligned}
\]
Durch Einsetzen der rechten Seite der rechten Seiten einer Gleichung
in die linke Seite der vorangehende Gleichung erhalten wir die
Darstellung
\begin{align*}
p(z)
&=
p^0(z)
\\
&=
p^1(z)(z- a) + A_1
\\
&=
p^2(z)(z- a)^2 + A_2(z- a) + A_1
\\
&=
p^3(z)(z- a)^3 + A_3 (z- a)^2 + A_2 (z- a) + A_1
\\[-2pt]
&\phantom{a}\vdots
\\
&=
A_m(z-a)^{m-1} + A_{m-1}(z- a)^{m-2} + \ldots + A_3 (z- a)^2 + A_2(z- a) + A_1.
\end{align*}
Damit kann jetzt auch die rationale Funktion einfacher ausgedrückt
werden:
\begin{align*}
f(z)
&=
\frac{p(z)}{(z- a)^n}
\\
&=
\frac{A_m}{(z- a)^{n-m+1}}
+
\frac{A_{m-1}}{(z- a)^{n-m+2}}
+
\ldots
+
\frac{A_3}{(z- a)^{n-2}}
+
\frac{A_2}{(z- a)^{n-1}}
+
\frac{A_1}{(z- a)^n}.
\end{align*}
Wir haben also gezeigt, dass eine rationale Funktion mit Nenner
$(z- a)^n$ immer geschrieben werden kann als eine Summe
von Termen der Form $A_k/(z- a)^{n-k}$.

%
% Potenz eines Polynoms im Nenner
%
\subsubsection{Potenz eines Polynoms im Nenner}
Der für den Nenner $q(z) = (z-a)^n$ durchgeführte Prozess lässt sich
auch für einen Nenner durchführen, der die Potenz $u(z)^n$ eines
einzigen Polynoms $u(z)$ ist.
Wir schreiben wieder $p^0(z)=p(z)$.
Der Divisionsalgorithmus zerlegt das Polynom $p^k(z)$ in 
\begin{equation}
p^k(z)
=
p^{k+1}(z)\cdot u(z) + A_{k+1}(z),
\label{buch:analytisch:rational:partial:eqn:u}
\end{equation}
wobei $\deg A_{k+1}(z)<\deg u(z)$ ist und
$\deg p^{k+1}(z) = \deg p^{k}(z)-\deg u(z)$.
Der Algorithmus bricht ab, sobald der Quotient $p^k(z)$ nicht
mehr weiter durch $u(z)$ geteilt werden kann, sobald also
sein Grad kleiner als der Grad von $u(z)$ ist.
Der Algorithmus bricht somit nach $m$ Schritten ab, wenn
$\deg p(z) - m\deg u(z) < \deg u(z)$ oder
$\deg p(z) < (m+1) \deg u(z)$.
Als Spezialfall beachten wir, dass $m=0$ ist, wenn der Nennergrad
$\deg p(z)$ bereits kleiner ist als der Grad von
$u(z)$: $\deg p(z) < \deg u(z)$.
Durch rekursives Einsetzen der Gleichungen
\eqref{buch:analytisch:rational:partial:eqn:u}
folgt
\begin{align*}
p(z)
&=
p^0(z)
\\
&=
p^1(z)\cdot u(z) + A_1(z)
\\
&=
\bigl(
p^2(z)\cdot u(z) + A_2(z)
\bigr)\cdot u(z) + A_1(z)
=
p^2(z)\cdot u(z)^2 + A_2(z) \cdot u(z) + A_1(z)
\\
&=
p^3(z)\cdot u(z)^2 + A_2(z)\cdot u(z) + A_1(z)
\\[-2pt]
&\phantom{a}\vdots
\\
&=
A_m(z)\cdot u(z)^{m-1} + A_{m-1}(z)\cdot u(z)^{m-2}
+
\ldots
+
A_3(z)\cdot u(z)^2 + A_2(z)\cdot u(z) + A_1(z).
\end{align*}
Durch Division durch $u(z)^n$ kann jetzt wieder ein Darstellung
\begin{align*}
f(z)
&=
\frac{p(z)}{u(z)^n}
\\
&=
\frac{A_m(z)}{u(z)^{n-m+1}}
+
\frac{A_{m-1}(z)}{u(z)^{n-m+2}}
+
\ldots
+
\frac{A_3(z)}{u(z)^{n-2}}
+
\frac{A_2(z)}{u(z)^{n-1}}
+
\frac{A_1(z)}{u(z)^n}
\end{align*}
gefunden werden, deren Nenner alle einen Grad $\deg A_k(z)<\deg u(z)$ 
kleiner als der Grad von $u(z)$ haben.

%
% Teilerfremde Faktoren im Nenner
%
\subsubsection{Teilerfremde Faktoren im Nenner}
Wir betrachten jetzt den Fall, dass der Nenner ein Produkt
$q(z) = q_1(z)q_2(z)$ von teilerfremden Faktoren $q_1(z)$ und $q_2(z)$
ist.
Der erweiterte euklidische Algorithmus zeigt, dass es Polynome
\index{euklidischer Algorithmus!erweitert}%
$s(z)$ und $t(z)$ gibt derart, dass
\begin{equation}
p(z)
=
s(z)\cdot q_1(z) + t(z)\cdot q_2(z)
\label{buch:analytisch:rational:partial:eukl}
\end{equation}
gilt.
Da $p(z)$ und $q(z)$ teilerfremd sind, kann $p(z)$ weder durch
$q_1(z)$ noch durch $q_2(z)$ teilbar sein.
Es folgt, dass auch
$s(z)$ nicht durch $q_2(z)$ teilbar sein kann, denn wäre $q_2(z)$
ein Teiler von $s(z)$, dann wäre $q_2(z)$ ein Teiler der rechten
Seite von
\eqref{buch:analytisch:rational:partial:eukl}
und damit auch von $p(z)$.

Mit der Darstellung
\eqref{buch:analytisch:rational:partial:eukl}
des Zählers ist die rationale Funktion
\[
f(z)
=
\frac{p(z)}{q(z)}
=
\frac{
s(z)\cdot q_1(z) + t(z)\cdot q_2(z)
}{q_1(z)\cdot q_2(z)}
=
\frac{s(z)}{q_2(z)}
+
\frac{t(z)}{q_1(z)}
\]
eine Summe von Brüchen, die nur die Faktoren $q_1(z)$ und $q_2(z)$
als Nenner haben.

Sei jetzt $q(z)=q_1(z)\cdot \ldots \cdot q_l(z)$ eine Faktorisierung
von $q(z)$ in paarweise teilerfremde Fatkoren.
Durch Wiederholung es oben dargestellten Prozesses kann die rationale
Funktion in die Form
\begin{equation}
f(z)
=
\frac{p(z)}{q(z)}
=
\frac{p_1(z)}{q_1(z)}
+
\ldots
+
\frac{p_l(z)}{q_l(z)}
\label{buch:analytisch:rational:partial:eqn:pksumqk}
\end{equation}
gebracht werden, wobei $p_k(z)$ nicht durch $q_k(z)$ teilbar ist.

%
% Allgemeine Partialbruchzerlegung
%
\subsubsection{Allgemeine Partialbruchzerlegung}
Im vorangegangenen Abschnitt wurde nicht verlangt, dass die Nenner
$q_k(z)$ irreduzibel sind.
Insbesondere ist für das Argument zulässig, dass jeder Nenner $q_k(z)$
ein Potenz $q_k(z) = u_k(z)^{n_k}$ ist.
Dann kann der Algorithmus des vorangegangenen Abschnitts auf jeden
Term auf der rechten Seite von
\eqref{buch:analytisch:rational:partial:eqn:pksumqk}
angewendet werden.
Es ergibt sich dann der folgende Satz.

\begin{satz}[Partialbruchzerlegung]
\label{buch:analytisch:rational:satz:partialbruchzerlegung}
Sei $f(z) = p(z)/q(z)$ eine rationale Funktion und 
\[
q(z)
=
q_1(z)^{n_1}
\cdot\ldots\cdot
q_l(z)^{n_l}
\]
eine Faktorisierung in teilerfremde Faktoren $q_k(z)$, $1\le k\le l$.
Dann gibt es ein Darstellung
\[
f(z)
=
p_0(z)
+
\sum_{k=1}^{m_1}
\frac{A_k^1(z)}{q_1(z)^{n_1-k+1}}
+
\ldots
+
\sum_{k=1}^{m_l}
\frac{A_k^l(z)}{q_l(z)^{n_1-k+1}}
\]
mit Polynomen $A_k^r(z)$
wobei $\deg A_k^r(z) < \deg q_k(z)$ ist.
\end{satz}

Man beachte, dass nirgends verlangt wird, dass die Faktoren 
$q_k(z)$ irreduzibel sein müssen.
Es wird nur verlangt, dass die Polynome keine gemeinsamen
Faktoren hat.
Dies ist mithilfe des euklidischen Algorithmus für Polynome
sehr leicht festzustellen.
Irreduzibilität ist dagegen algorithmisch viel schwieriger
zu etablieren.

Im Integrationsproblem ermöglicht die Partialbruchzerlegung
wegen der Linearität des Integrals, die Aufgabe in separate
Integrationsprobleme für rationale Funktionen mit
einfacheren Nennern zu zerlegen.
Die klassische Methode, die man in der Analysis lernt,
geht davon aus, dass sich dank des Fundamentalsatzes der Algebra
das Nennerpolynom in Linearfaktoren zerlegen lässt, so dass alle
$q_k(z)$ von der Form $z-a_k$ sind.
Es müssen dann nur noch Integrale der Form
\[
\int
\frac{A}{(z-a)^n}
\,dz
=
\begin{cases}
\displaystyle-\frac{1}{n-1} \frac{A}{(z-a)^{n-1}}&\qquad\text{für $n>1$} \\
A\log (z-a)                                      &\qquad\text{für $n=1$}
\end{cases}
\]
bestimmt werden.
Die Partialbruchzerlegung löst also mindestens im Prinzip
das Integrationsproblem für rationale Funktionen.
Leider gibt es für Polynome vom Grad grösser als $5$ im
Allgemeinen keine Formel zur Bestimmung der Nullstellen.
Somit gibt es für ein Computeralgebrasystem im Allgmeinen
keine Möglichkeit, ein Integral auf diesem Weg zu bestimmen.

%
% Partialbruchzerlegung und lineare Gleichungssysteme
%
\subsubsection{Partialbruchzerlegung und lineare Gleichungssysteme}
Die in den vorangegangenen Abschnitten beschriebene Methode zur
Bestimmung der Partialbruchzerlegung verlangt die wiederholte
Durchführung des euklidischen Algorithmus für Polynome.
Da die Form der Partialbruchzerlegung bereits bekannt ist,
kann sie mit unbekannten Koeffizienten als
\begin{equation}
\frac{p(z)}{(z-a_1)^{n_1}\cdots(z-a_l)^{n_l}}
=
\sum_{k=1}^l
\sum_{s=1}^{n_l}
\frac{A^k_s}{(z-a_k)^s}
\label{buch:analytisch:rational:partial:eqn:ansatz}
\end{equation}
angesetzt werden.
Man beachte, dass die Koeffizienten $A^k_s$ in
\eqref{buch:analytisch:rational:partial:eqn:ansatz} nicht
die gleiche Bedeutung haben wie in der Herleitung der
Partialbruchzerlegung in den vorangegangenen Abschnitten.
Durchmultiplizieren mit dem Nenner $(z-a_1)^{n_1}\dots (z-a_l)^{n_l}$
und Koeffizientenvergleich liefert dann ein lineare Gleichungssystem
für die Koeffizienten $A^k_s$.
\index{lineares Gleichungssystem}%

\begin{beispiel}
Wir bestimmen die Partialbruchzerlegung der rationalen Funktion
\begin{equation}
f(z)
=
\frac{(z+2)(z-2)}{(z-1)^2(z+1)^3}
=
\frac{z^2-4}{z^5+z^4-2z^3-2z^2+z+1}.
\label{buch:analytisch:rational:partial:bsp:aufgabe}
\end{equation}
Dazu verwenden wir nach dem Muster von
\eqref{buch:analytisch:rational:partial:eqn:ansatz}
den Ansatz
\begin{equation}
f(z)
=
\frac{A^1_1}{z-1}
+
\frac{A^1_2}{(z-1)^2}
+
\frac{A^2_1}{z+1}
+
\frac{A^2_2}{(z+1)^2}
+
\frac{A^2_3}{(z+1)^3},
\label{buch:analytisch:rational:partial:bsp:ansatz}
\end{equation}
multiplizieren mit $q(z)$ und führen den Koeffizientenvergleich
mit den Koeffizienten des Polynoms $p(z)=z^2-4$ durch.
Die etwas längere Rechnung kann mit einem Computeralgebrasystem
vereinfacht werden und liefert das Gleichungssystem
\[
\renewcommand{\arraycolsep}{2pt}
\renewcommand{\arraystretch}{1.3}
\begin{array}{rcrcrcrcrcr}
  A^1_1 & &        &+&  A^2_1 & &       & &        &=&  0\\
 2A^1_1 &+&  A^1_2 & &        &+& A^2_2 & &        &=&  0\\
        & & 3A^1_2 &-& 2A^2_1 &-& A^2_2 &+&  A^2_3 &=&  1\\
-2A^1_1 &+& 3A^1_2 & &        &-& A^2_2 &-& 2A^2_3 &=&  0\\
- A^1_1 &+&  A^1_2 &+&  A^2_1 &+& A^2_2 &+&  A^2_3 &=& -4\\
\end{array}
\]
Die Lösung des Gleichungssystems ergibt
\[
A^1_1 =  \frac{13}{16},\quad
A^1_2 = -\frac{ 3}{ 8},\quad
A^2_1 = -\frac{13}{16},\quad
A^2_2 = -\frac{ 5}{ 4},\quad
A^2_3 = -\frac{ 3}{ 4}.
\]
Setzt man diese Koeffizienten in den Ansatz
\eqref{buch:analytisch:rational:partial:bsp:ansatz}
ein, erhält man
\[
f(z)
=
\frac{13}{16(z-1)  }
-
\frac{ 3}{ 8(z-1)^2}
-
\frac{13}{16(z+1)  }
-
\frac{ 5}{ 4(z+1)^2}
-
\frac{ 3}{ 4(z+1)^3}.
\]
Gleichnamig machen liefert wieder den ursprünglichen Bruch
\eqref{buch:analytisch:rational:partial:bsp:aufgabe}
\end{beispiel}

